% ornek.tex -- gtu-el-yazisi şablonu için 2 sayfalık örnek
% Derleme: latex-sablon/derle.sh latex-sablon/ornek.tex
\documentclass[11pt,a4paper]{article}
\usepackage{gtu-el-yazisi}

\begin{document}

\notbilgi[2023--24 Güz $\cdot$ örnek sayfa]{Devre Analizi I}{6. Hafta -- RC Devreleri ve Fazörler}{devre1_hafta6_notlar.pdf}

\section{Birinci Dereceden RC Devresi}
\sayfa{1}
\begin{minipage}[c]{0.52\linewidth}
  Anahtar $t=0$ anında kapatılıyor; kondansatör başlangıçta boş, yani $v_C(0^-)=0$.
  \vurgu{Kondansatör gerilimi ani olarak değişemez}, bu yüzden $v_C(0^+)=v_C(0^-)$.
\end{minipage}\hfill
\begin{minipage}[c]{0.44\linewidth}
  \centering
  \begin{circuitikz}[line width=0.7pt]
    \draw (0,0) to[V, l=$V_s$] (0,2.2)
          to[cspst, l={$t=0$}] (2.2,2.2)
          to[R, l=$R$, i>^=$i(t)$] (4.4,2.2)
          to[C, l_=$C$, v^=$v_C(t)$] (4.4,0) -- (0,0);
  \end{circuitikz}
\end{minipage}

KGY (Kirchhoff Gerilim Yasası) çevreye uygulanırsa:
\begin{align}
  V_s &= R\,i(t) + v_C(t), \qquad i(t) = C\,\frac{dv_C}{dt} \\
  \Rightarrow\quad RC\,\frac{dv_C}{dt} + v_C &= V_s \qquad (t \ge 0)
\end{align}

\begin{tanim}[(Zaman sabiti)]
  $\tau = RC$ devrenin \textbf{zaman sabitidir}; birimi saniyedir:
  $[\si{\ohm}]\cdot[\si{\farad}] = [\si{\second}]$. Büyük $\tau$ $\Rightarrow$ yavaş şarj.
\end{tanim}

\formulkutu[ezberle!]{v_C(t) = V_s\left(1 - e^{-t/\tau}\right), \qquad t \ge 0}

\begin{ornek}[1]
  $R = \qty{10}{\kilo\ohm}$, $C = \qty{4,7}{\micro\farad}$ ve $V_s = \qty{12}{\volt}$ ise
  $t = \qty{47}{\milli\second}$ anında $v_C$ kaç volttur?
\end{ornek}
\begin{cozum}
  $\tau = RC = (10\cdot10^{3})(4{,}7\cdot10^{-6}) = \qty{47}{\milli\second}$, yani $t = \tau$:
  \[ v_C(\tau) = 12\,\bigl(1 - e^{-1}\bigr) \approx 12 \times 0{,}632
     = \kutula{\qty{7,59}{\volt}} \]
\end{cozum}

\noindent
\begin{minipage}[c]{0.36\linewidth}
  \centering
  \begin{tabular}{cc}
    \toprule
    $t$ & $v_C / V_s$ \\ \midrule
    $\tau$  & \%63,2 \\
    $2\tau$ & \%86,5 \\
    $3\tau$ & \%95,0 \\
    $5\tau$ & \%99,3 \\ \bottomrule
  \end{tabular}
  \captionof{table}{Şarj yüzdeleri}
\end{minipage}\hfill
\begin{minipage}[c]{0.6\linewidth}
  \centering
  \begin{tikzpicture}
    \begin{axis}[width=7.4cm, height=4cm, xmin=0, xmax=5.4, ymin=0, ymax=1.18,
        xlabel={$t/\tau$}, ylabel={$v_C/V_s$}, xtick={1,2,3,4,5},
        ytick={0.632,1}, yticklabels={$0{,}632$,$1$}]
      \addplot[domain=0:5.2, samples=60, smooth] {1-exp(-x)};
      \addplot[dashed, kursunkalem, line width=0.6pt] coordinates {(0,0.632) (1,0.632) (1,0)};
    \end{axis}
  \end{tikzpicture}
\end{minipage}

\begin{dikkat}
  $5\tau$ sonra kondansatör ``dolmuş'' kabul edilir, ama matematiksel olarak
  $v_C$ hiçbir zaman tam $V_s$ olmaz (asimptot). Sınavda \vurgu[pembe]{başlangıç koşulunu
  $v_C(0^+)$ yazmayı unutma!}
\end{dikkat}

\section{Fazörler ve Empedans}
\sayfa{2}
Sinüzoidal kararlı halde her işaret bir fazörle gösterilir:
$v(t) = V_m\cos(\omega t + \phi) \donusum \mathbf{V} = V_m\phase{\phi}$, burada $\omega = 2\pi f$.
Empedanslar: $Z_R = R$, \ $Z_L = j\omega L$, \ $Z_C = \dfrac{1}{j\omega C} = -\dfrac{j}{\omega C}$.

\begin{teorem}[(Thevenin)]
  Doğrusal iki uçlu her devre, bir $\mathbf{V}_{th}$ kaynağı ile seri bir $Z_{th}$ empedansına
  eşdeğerdir: $\mathbf{V}_{th}$ açık devre gerilimi, $Z_{th} = \mathbf{V}_{th}/\mathbf{I}_{kd}$.
\end{teorem}

Düğüm gerilimleri yöntemi (matris biçimi), $\alpha$, $\beta$, $\theta$ gibi açılar derece cinsinden:
\[
  \begin{bmatrix}
    \dfrac{1}{R_1} + \dfrac{1}{R_2} & -\dfrac{1}{R_2} \\[2mm]
    -\dfrac{1}{R_2} & \dfrac{1}{R_2} + j\omega C
  \end{bmatrix}
  \begin{bmatrix} \mathbf{V}_1 \\ \mathbf{V}_2 \end{bmatrix}
  =
  \begin{bmatrix} \mathbf{I}_s \\ 0 \end{bmatrix},
  \qquad
  \Delta = \det(\mathbf{Y}) \neq 0
\]

\sekilnotu[3]{Hocanın tahtaya çizdiği fazör diyagramı: $\mathbf{V}_R$, $\mathbf{V}_L$ ve
$\mathbf{V}_C$ vektörleri, aralarında $90^\circ$ faz farkı}

\subsection{Laplace ile bağlantı}
$s = j\omega$ konursa empedanslar Laplace düzlemine taşınır: $Z_L(s) = sL$, $Z_C(s) = 1/(sC)$.
\[
  \Lap{e^{-at}u(t)} = \int_{0}^{\infty} e^{-at} e^{-st}\,dt = \frac{1}{s+a}, \qquad
  \Fou{\delta(t)} = 1, \qquad \sum_{n=0}^{\infty} x^n = \frac{1}{1-x}\ (|x|<1)
\]
Hoca burada bir şey daha ekledi: \okunamadi{rezonans frekansı} için
$\omega_0 = 1/\sqrt{LC}$ ve kalite faktörü $Q = \omega_0 L / R$ \okunamadi{}.

\textit{English note from the lecture:} ``The phasor method is valid only in sinusoidal steady
state; transients must be handled in the time domain or with the Laplace transform.''

\begin{ozet}
  \begin{itemize}
    \item RC şarj: $v_C(t) = V_s(1-e^{-t/\tau})$, \ $\tau = RC$, \ $5\tau$'da $\approx$ \%99.
    \item Fazör: $\cos(\omega t+\phi) \to 1\phase{\phi}$; türev $\to$ $j\omega$ ile çarpma.
    \item $Z_L = j\omega L$ (gerilim akımdan $90^\circ$ önde), $Z_C = 1/(j\omega C)$ ($90^\circ$ geride).
  \end{itemize}
\end{ozet}

\subsection*{Orijinal el çizimi}
El ile çizilmiş şekiller orijinal sayfadan kırpılıp aynen eklenir:
\elcizimi[0.55]{ornek-cizim.png}[3]

\end{document}
